\renewcommand{\theposttitle}{What’s the difference between generosity and strategy?}
\renewcommand{\thepostdate}{2025-12-16}
\vspace*{\fill}
\begin{center}
\begin{minipage}{0.92\textwidth}
\small
The best visual metaphors reveal an insight that \textbf{feels new, but still oddly familiar. }They remind us of something we quietly know in our mind, heart, or soul, but due to the frantic hustle \& bustle of the day to day, we forget!
\end{minipage}
\end{center}
\vspace{4mm}
\begin{center}
\includegraphics[width=\textwidth, height=0.24\textheight, keepaspectratio]{C:/Users/James/Documents/GitHub/VisualIdeasBook/data/images/45f3c04dbcf151c6bdd7a225e2fa9cc0.png}
\end{center}
\vspace{2mm}
\begin{center}
\begin{minipage}{0.92\textwidth}
\small
The key is using what’s familiar to act as the relational bridge between what we know and what feels foreign. Usually the bridge connects (or corrects) a behavior or belief.
\end{minipage}
\end{center}
\vspace{5mm}
\begin{center}
\includegraphics[width=\textwidth, height=0.24\textheight, keepaspectratio]{C:/Users/James/Documents/GitHub/VisualIdeasBook/data/images/8e9040af525f88582daecec111a347a2.png}
\end{center}
\vspace{2mm}
\begin{center}
\begin{minipage}{0.92\textwidth}
\small
For example, if you’re struggling to find pockets of happiness in your life and I just said “fix your environment and you might be a little happier,” it doesn’t feel very insightful.

So instead, I start with an assumption of what you might familiar (and it’s helpful to point out no assumption holds true for everyone). But what if I assumed you might know what it feels like to chase a butterfly?
\end{minipage}
\end{center}
\vspace*{\fill}